\section{Setup and identification geometry}
\label{sec:setup}

Let $Z_t=Y_{t+1}-Y_t$ and let all conditional statements hold for $P_X$-almost every $x$.

\begin{assumption}[Pooled conditional-mean model]
\label{ass:mean}
The state satisfies $Y_t\in[0,1]$. There exist measurable bounded functions $U_0,R_0:\cX\to[0,\infty)$ such that
\begin{equation}
\E[Z_t\mid X_t=x,Y_t=y]=U_0(x)(1-y)-R_0(x)y,
\label{eq:mean_model}
\end{equation}
and the residual has conditional mean zero and finite conditional variance.
\end{assumption}

Assumption~\ref{ass:mean} is predictive, not causal. When observations from multiple counties are pooled, it additionally requires that the two conditional-mean functions are common after conditioning on $X$, or that $X$ contains sufficient unit-level modifiers. If unobserved county composition changes with $Y$, the pooled conditional mean need not remain affine in $y$; the fitted components are then projections rather than physical rates.

At a fixed $x$, the two-flow class and the signed one-flow class are
\begin{align}
\cF_2(x)&=\{u(1-y)-ry:u,r\ge0\},\label{eq:f2}\\
\cF_1(x)&=\{a(1-y):a\ge0\}\cup\{-by:b\ge0\}.\label{eq:f1}
\end{align}
Define
\begin{align}
A(x)&=\E[(1-Y)^2\mid X=x],\nonumber\\
B(x)&=\E[Y^2\mid X=x],\nonumber\\
C(x)&=\E[Y(1-Y)\mid X=x],\nonumber\\
v(x)&=\Var(Y\mid X=x).
\label{eq:moments}
\end{align}
Writing $\mu(x)=\E[Y\mid X=x]$ gives $A=1-2\mu+\E[Y^2]$, $B=\E[Y^2]$, and $C=\mu-\E[Y^2]$, hence
\begin{equation}
A(x)B(x)-C(x)^2=v(x).
\label{eq:moment_identity}
\end{equation}

The natural coordinates are net drift at the conditional mean state and total turnover,
\begin{align}
\alpha(x)&=(1-\mu(x))U_0(x)-\mu(x)R_0(x),\nonumber\\
\tau(x)&=U_0(x)+R_0(x).
\label{eq:alpha_tau}
\end{align}
Then
\begin{equation}
m_0(x,y)=\alpha(x)+\tau(x)\{\mu(x)-y\}.
\label{eq:orthogonal_mean}
\end{equation}
The inverse transformation is $U_0=\alpha+\mu\tau$ and $R_0=(1-\mu)\tau-\alpha$.

\begin{proposition}[Orthogonal identification geometry]
\label{prop:geometry}
Let $\psi_x(Y)=(1,\mu(x)-Y)^\top$. Then
\begin{equation}
\E[\psi_x(Y)\psi_x(Y)^\top\mid X=x]
=\begin{bmatrix}1&0\\0&v(x)\end{bmatrix}.
\label{eq:orthogonal_information}
\end{equation}
Equivalently, for $\varphi(Y)=(1-Y,-Y)^\top$ and $Q(x)=\E[\varphi(Y)\varphi(Y)^\top\mid X=x]$,
\begin{align}
\det Q(x)&=v(x),\nonumber\\
v(x)&\le\lambda_{\min}(Q(x))\le2v(x).
\label{eq:eigen_bounds}
\end{align}
with $1/2\le\lambda_{\max}(Q(x))\le1$. Consequently, the two conditional-mean components are separately identified if and only if $v(x)>0$.
\end{proposition}

The proposition separates a well-identified net drift from a potentially weak turnover direction. In a centered homoskedastic fixed design with $n$ observations,
\begin{align}
\Var(\widehat\alpha\mid Y)&=\frac{\sigma_\varepsilon^2}{n},\nonumber\\
\Var(\widehat\tau\mid Y)&=\frac{\sigma_\varepsilon^2}{n\widehat v},\nonumber\\
\Cov(\widehat\alpha,\widehat\tau\mid Y)&=0.
\label{eq:alpha_tau_variance}
\end{align}
Thus accurate prediction of net drift can coexist with unstable recovery of interruption and restoration. Proposition~\ref{prop:geometry} is a population identification statement; continuous drivers require function sharing across nearby $x$, and Equation~\eqref{eq:alpha_tau_variance} is a fixed-design benchmark rather than a finite-sample theorem for the dynamic neural panel.
